\subsection{The actual signed source and its Fermat actions}

We now use the actual automorphism source of \S49.IV--V, rather
than the minimal parameter scheme of \S51.2. We inspected the
original images at cote 140--3, scans 530, 534, 549, 550, 559,
and 560. Their numbered folios are 568, 570, 578, and 583;
the intervening unnumbered versos are recorded without silently
renumbering them. The definition of the linked signs is (7) on
scan 530; the enlarged numerical source is (6) on scan 549;
the fibre products are (50)--(59) on scans 559--560
\citep[\S49.IV--V]{GrothendieckLM1403Original}.
As in B.25, a scan number means the physical PDF page; its
archive footer is one smaller.

Write $C=\widehat{\mathrm{SL}_2(\mathbb Z)}$ and use
\[
 \rho^6=\sigma^4=1,\qquad
 \omega=\rho^3=\sigma^2,\qquad
 e_0=\sigma\rho,\quad e_1=\rho\sigma,
 \quad x=\omega e_0^2,\quad y=\omega e_1^2.
\]
Let $F=\overline{\langle x,y\rangle}$, the free profinite
rank-two normal subgroup of index twelve from B.24a, and set
\[
 K_n=\ker\bigl(F\longrightarrow(\mathbb Z/n)^2\bigr)
       =\overline{F'F^n},\qquad H_n=C/K_n.
\]
Every derived subgroup in this subsection is closed.
For $\mu\in\widehat\mathbb Z^\times$ let $\epsilon_3(\mu)$
and $\epsilon_2(\mu)$ be its signs modulo three and four.
We retain the sign order $(\rho,\sigma)$, which reverses
the order used in some manuscript displays.

\begin{statement}{B.31}{Proved}{Effective source, preserved kernels, and continuous relative actions}
The two actual sources
\[
\begin{aligned}
D_B={}&\{(u,\mu):u\in\operatorname{Aut}_{\rm cts}(C),\
 u(e_0)=e_0^\mu,\
 u(\rho)\sim\rho^{\epsilon_3(\mu)},\
 u(\sigma)\sim\sigma^{\epsilon_2(\mu)}\},\\
D_Z={}&\{(u,\mu,e_r,e_s):u\in\operatorname{Aut}_{\rm cts}(C),\
 u(e_0)=e_0^\mu,\
 u(\rho)\sim\rho^{e_r},\
 u(\sigma)\sim\sigma^{e_s}\}
\end{aligned}
\]
are naturally identical for this $C$: the two signs of $D_Z$
are forced by $\mu$. The numerical coordinates are determined
by $u$. Denote the resulting profinite group by $D$.
Every $u\in D$ fixes $\omega$, sends $x$ to $x^\mu$, and
preserves $F$, every $K_n$, and every closed characteristic
subgroup of $K_n$.

Consequently it acts functorially on
\[
 C/[K_n,K_n],\qquad
 R_{n,\ell}=C/\ker(K_n\longrightarrow K_n^{(\ell)}),
 \qquad \mathcal R_{n,\ell}
                                                        \tag{B.31.1}
\]
for all $n\ge1$ and primes $\ell$. Here $R_{n,\ell}$ is
a profinite group with free pro-$\ell$ kernel of rank $n^2+1$
over the finite quotient $H_n$, and $\mathcal R_{n,\ell}$ is
its continuous relative prounipotent completion over
$\mathbb Q_\ell$ with the same constant quotient.
Its unipotent radical is free prounipotent of rank $n^2+1$.
The action also extends to the canonical abelian-kernel
group-scheme stage over $\mathbb Z_\ell$.
These are actions on the actual relative objects, with
their induced finite-quotient actions, at every level.

Let $L_0=\overline{\langle e_0\rangle}$, let
$j(\ell)=\operatorname{Int}(\ell)\in D$, and put
\[
 \gamma=D/j(L_0),\qquad
 G_D=(C\rtimes D)/\{(\ell^{-1},j(\ell)):\ell\in L_0\}.
                                                        \tag{B.31.2}
\]
There is a specified action
$\theta:G_D\to\operatorname{Aut}_{\rm cts}(C)$,
$[c,u]\mapsto\operatorname{Int}(c)u$, which induces all
the actions (B.31.1). The signed normalizers and the
fibre products of scans 559--560 inherit actions through
their stated projections to $G_D$; the common quotient
$\gamma$ has a canonical outer action. This assertion
retains the conjugators and the choice of strict action.
It does not identify these sources with the minimal
\S51 parameter scheme.
\end{statement}
\begin{proof}
The abelianization of $C$ is $\mathbb Z/12$; one can take
$[\rho]=2$, $[\sigma]=-3$, and then $[e_0]=-1$ is a
generator. Abelianizing $u(e_0)=e_0^\mu$ shows that $u$
acts by $\mu$ on this cyclic group. The two conjugacy
conditions therefore imply
\[
        e_r\equiv\mu\pmod6,\qquad e_s\equiv\mu\pmod4.
\]
They force precisely the two signs in $D_B$.
Moreover $u(\omega)=u(\sigma)^2=\omega$ since $\omega$
is central and $e_s$ is odd. Because $\mu$ is odd,
\[
             u(x)=\omega e_0^{2\mu}
                         =(\omega e_0^2)^\mu=x^\mu.
                                                        \tag{B.31.3}
\]

The group $F$ is the closed normal closure of $x$ in $C$:
it is normal and generated by $x,y$, and $y=\sigma x\sigma^{-1}$.
The closed cyclic groups generated by $x$ and $x^\mu$ coincide.
Thus (B.31.3) proves $u(F)=F$, and the characteristic formula
for $K_n$ proves preservation of every stated kernel.
In particular this does not assume preservation before
applying a completion's universal property.

The procyclic group $L_0$ is isomorphic to
$\widehat\mathbb Z$. Indeed the relation
$e_0^{2a}=\omega^a x^a$ and the primitive class of $x$ in
$F^{\rm ab}=\widehat\mathbb Z^2$ show that $e_0^a=1$
implies $a=0$. Thus $u(e_0)$ determines $\mu$.
The equations defining $D$ are closed equations and
closed conjugacy conditions in the profinite automorphism
group of the finitely generated group $C$, so $D$ has
the claimed profinite topology. The same argument works
if the alternative loop $\sigma^{-1}\rho=\omega e_0$
is used: its square is $e_0^2$, and all multipliers are odd.
No doubtful sign in scan 530 changes this conclusion.

The quotient $H_n$ has order $12n^2$, while Schreier's
formula gives $\operatorname{rank}K_n=n^2+1$. A finite-index
subgroup of the discrete free group induces its full
profinite topology, so $K_n$ is the free profinite group
of this rank. Its pro-$\ell$ and continuous prounipotent
completions are consequently free on the same generators.
Their defining kernels are characteristic in $K_n$.
This proves the profinite actions in (B.31.1).

We give a concrete construction of the algebraic action.
Choose coset representatives $t_h$ for $C/K_n$. Then
conjugation by $t_h$ induces $a_h$ on $K_n$, and
$c(h,k)=t_ht_kt_{hk}^{-1}\in K_n$. Push both data through
$K_n\to U_\ell(K_n)$, where $U_\ell$ denotes continuous
prounipotent completion. The disjoint union of copies of
$U_\ell(K_n)$ indexed by $H_n$, with law
\[
       (v,h)(w,k)=(v\,a_h(w)c(h,k),hk),
\]
is the required proalgebraic group $\mathcal R_{n,\ell}$.
The usual crossed-action and cocycle identities are
identities of the source group and therefore still hold.
They also prove its relative universal property.
For $u\in D$, write
$u(t_h)=b_h t_{\bar u(h)}$, with $b_h\in K_n$.
The formula
\[
             (v,h)\longmapsto(u_U(v)b_h,\bar u(h))
                                                        \tag{B.31.4}
\]
defines the action; the inverse is obtained from $u^{-1}$.
All maps on $U_\ell$ exist by the continuous free universal
property. In particular infinite profinite words are
evaluated continuously over $\mathbb Q_\ell$, not as
unspecified rational words.
At every finite nilpotent level, the images of a finite
free basis depend continuously on $u$, and the ensuing
group operations are $\ell$-adic polynomial operations.
This proves continuity of the action in the inverse-limit
topology on these levels.

For the abelian stage, $K_n^{\rm ab}=\widehat\mathbb Z^{n^2+1}$.
Project its action matrices, translations $b_h$, and
integral coset cocycle to their $\ell$ components.
Formula (B.31.4) is now an affine group-scheme
automorphism over $\mathbb Z_\ell$. Its identities and
inverse hold coefficientwise. This argument is confined
to the indicated integral abelian stage; it does not
replace higher integral groups by a BCH lattice at a
prime dividing the BCH denominators.

For completeness we type the retained conjugator sources.
The intersection $L_0\cap Z(C)$ is trivial. To see this,
work already in $C/F'$. Conjugation by $e_0$ on
$F^{\rm ab}$ has matrix
$J=\left(\begin{smallmatrix}1&-1\\0&-1\end{smallmatrix}\right)$.
A central $e_0^a$ has even $a$, so write $a=2b$.
Then $e_0^{2b}=\omega^b x^b$ is central only if
$(1-R)b(1,0)^t=0$, where
$R=\left(\begin{smallmatrix}0&-1\\1&-1\end{smallmatrix}\right)$.
This forces $b=0$. Thus $C$ embeds in $G_D$, and
\[
                1\longrightarrow C\longrightarrow G_D
                    \longrightarrow\gamma\longrightarrow1
\]
is exact. The antidiagonal in (B.31.2) is closed and
normal, and its action $\operatorname{Int}(\ell^{-1})
\operatorname{Int}(\ell)$ is the identity. Hence $\theta$
is well-defined. Inner automorphisms also preserve every
$K_n$ and every kernel already used, proving its actions.

For $z=\rho,\sigma$, define
\[
 N_z^*=\{(g,e):g\in G_D,\ e=\epsilon_z(\chi(g)),\
                    \theta(g)(z)=z^e\},
\]
where $\epsilon_\rho=\epsilon_3$, $\epsilon_\sigma=\epsilon_2$,
and $\chi:G_D\to\widehat\mathbb Z^\times$ is the multiplier.
The signs and multiplier descend through $\gamma$.
The groups and maps relevant to the manuscript are
\[
\begin{aligned}
 \Gamma_D&=N_\rho^*\times_\gamma N_\sigma^*,&
 S_0\Gamma_D&=D\times_\gamma\Gamma_D,\\
 \mathcal M_D&=\Gamma_D\times_\gamma G_D,&
 S_0\mathcal M_D&=D\times_\gamma\mathcal M_D.
\end{aligned}
\]
The normalizers act by their $G_D$ projection;
$\mathcal M_D$ and $S_0\mathcal M_D$ act by the displayed
last $G_D$ coordinate. The $D$ coordinate of
$S_0\Gamma_D$ gives its based action. The two projections
of $\Gamma_D$ give two specified actions differing by
an inner automorphism. Thus its outer action, and that
of $\gamma$, are unambiguous. These assertions follow
from the explicit maps, without discarding any numerical
or conjugator coordinate and without an automatic
realization claim for a proposed pair of generator images.
\end{proof}

\begin{statement}{B.32}{Proved}{What the two full Fermat families detect}
The whole profinite tower of abelian-kernel Fermat stages is
exactly the relative metabelian quotient:
\[
       \bigcap_{n\ge1}[K_n,K_n]=F'',\qquad
       C/F''\ \simeq\ \varprojlim_{n\mid m} C/[K_n,K_n].
                                                        \tag{B.32.1}
\]
This is an equivariant isomorphism for the sources in B.31.
For example its action kernel in $D$ is exactly
\[
 \{u\in D:u(\rho)\rho^{-1}\in F'',\
                     u(\sigma)\sigma^{-1}\in F''\}.
                                                        \tag{B.32.2}
\]

The collection of all relative pro-$\ell$ Fermat quotients,
with both $n$ and $\ell$ varying, detects exactly the finite
quotients of $F$ whose derived subgroup is nilpotent.
Equivalently, if $\mathcal N\mathcal A$ denotes this class
of finite groups and $F^{\mathcal N\mathcal A}$ its maximal
pro-$\mathcal N\mathcal A$ quotient, then
\[
 \ker(F\to F^{\mathcal N\mathcal A})
       =\bigcap_{n,\ell}\ker(F\to R_{n,\ell}).
                                                        \tag{B.32.3}
\]
This quotient is proper: the full free profinite group is
not recovered by either collection. Formula (B.32.2)
is a concrete word-congruence criterion, not a classification
of all liftable automorphisms of $C/F''$.
\end{statement}
\begin{proof}
Since $F'\subset K_n$, the inclusion
$F''\subset[K_n,K_n]$ holds for every $n$. Conversely,
let $g\notin F''$. The case $g\notin F'$ is detected
already at $n=1$. Otherwise choose a finite quotient
$Q$ of the profinite metabelian group $F/F''$ in which
$g$ has nontrivial image. Choose $n$ divisible by the
exponent of $Q^{\rm ab}$. The subgroup $K_n$ maps into
$Q'$, so $[K_n,K_n]$ maps into $Q''=1$. This excludes
$g$ from the intersection and proves the first equality.

The natural map in (B.32.1) is therefore injective.
For a compatible tuple of quotient elements, its
preimages in the compact group $C$ are closed nonempty
cosets with the finite-intersection property: a common
multiple of finitely many levels suffices. Their
intersection is nonempty, proving surjectivity.
The action is equivariant because every subgroup in
the construction was proved invariant in B.31.
The two generators $\rho,\sigma$ topologically generate
the quotient, which proves (B.32.2).

For (B.32.3), the image of $F$ in $R_{n,\ell}$ is
an extension of $(\mathbb Z/n)^2$ by a pro-$\ell$
group. Its finite quotients have derived subgroup an
$\ell$-group, and therefore belong to
$\mathcal N\mathcal A$. This class is closed under
subgroups, quotients, and finite products. In particular
all finite quotients detected by finite collections
of these relative quotients remain in the class.

Conversely let $F\twoheadrightarrow Q$ be any finite
quotient with $Q'$ nilpotent. For a prime $\ell$ dividing
$|Q'|$, let $T_\ell$ be the product of its Sylow
subgroups at primes other than $\ell$. It is characteristic
in $Q'$ and normal in $Q$. The quotient $Q/T_\ell$ has
derived subgroup an $\ell$-group. Choose $n$ divisible
by the exponent of $Q^{\rm ab}$. The image of $K_n$
in $Q/T_\ell$ is an $\ell$-group, so the homomorphism
factors through the corresponding relative pro-$\ell$
quotient. The intersection of the $T_\ell$ is trivial.
If $Q'=1$, any prime and a multiple of the exponent
of $Q$ suffice. Thus every finite quotient in the
class is jointly detected, proving (B.32.3).

Finally $F$ maps onto $A_5$: take the five-cycle
$(12345)$ and the three-cycle $(123)$. The exact
permutation computation in \texttt{b\_signed\_source.py}
gives order sixty, derived subgroup of order sixty,
and trivial centre. Hence its derived subgroup is
not nilpotent. A pro-$\mathcal N\mathcal A$ group
cannot have this finite quotient. Were the kernel
in (B.32.3) trivial, compactness would identify $F$
with its maximal pro-$\mathcal N\mathcal A$ quotient,
a contradiction. This proves properness without
mistaking absence of one isolated quotient map for
a statement about an arbitrary injective dense map.
\end{proof}

\begin{statement}{B.33}{Proved}{The exact actual-source image on the first modular stage}
Put $\Delta=C/F'$. There are canonical automorphisms
$s_\mu$ of $\Delta$, for every
$\mu\in\widehat\mathbb Z^\times$, characterized by
\[
       s_\mu(a)=a^\mu,\qquad s_\mu(b)=b^\mu,
       \qquad a=e_0^{-1},\quad b=e_1^{-1}.
                                                        \tag{B.33.1}
\]
They satisfy $s_\mu s_\nu=s_{\mu\nu}$.
The image of the full actual signed source is exactly
\[
 \operatorname{im}(D\to\operatorname{Aut}(\Delta))
  =\{\operatorname{Int}(e_0^r)s_\mu:
          r\in\widehat\mathbb Z,\ \mu\in\widehat\mathbb Z^\times\}
  \simeq\widehat\mathbb Z\rtimes\widehat\mathbb Z^\times,
                                                        \tag{B.33.2}
\]
with law $(r,\mu)(s,\nu)=(r+\mu s,\mu\nu)$.
The pair $(r,\mu)$ is unique. Its kernel in $D$ is
the explicitly based congruence subgroup
\[
 D_0=\{u\in D:u(\rho)\rho^{-1},\
                       u(\sigma)\sigma^{-1}\in F'\}.
                                                        \tag{B.33.3}
\]
This determines the image on the first modular stage;
it does not determine the retained conjugator images
on the first Fermat stage or the full higher tower.
\end{statement}
\begin{proof}
In the braid notation of E.2 and E.6,
$\rho=ab$, $\sigma=(aba)^{-1}$, and the braid relation
gives $e_0=a^{-1}$, $e_1=b^{-1}$. The central braid
$z=(ab)^3$ maps to $\omega$, and $C$ is obtained by
the relation $z^2=1$. The standard profinite GT action is
\[
 a\longmapsto a^\mu,\qquad
 b\longmapsto f(a^2,b^2)^{-1}b^\mu f(a^2,b^2),
 \qquad f\in\widehat F_2'.
\]
Its multiplier map contains the cyclotomic character
of $G_\mathbb Q$, which is onto
$\widehat\mathbb Z^\times$: for every $m$ the finite
cyclotomic extension has Galois group
$(\mathbb Z/m)^\times$, and compactness gives
surjectivity on the inverse limit. Thus every $\mu$
occurs. The action sends $z$ to $z^\mu$, so it
descends to $C$; E.6 supplies the finite-factor
conjugacy conditions. It fixes $e_0$ to its indicated
power and therefore belongs to $D$.

In $\Delta$, the elements
$a^2=\omega x^{-1}$ and $b^2=\omega y^{-1}$ commute.
Consequently $f(a^2,b^2)=1$. The induced action
depends only on $\mu$ and is precisely (B.33.1).
These formulas on the two topological generators
also give $s_\mu s_\nu=s_{\mu\nu}$ and
$s_\mu(e_0)=e_0^\mu$.

It remains to classify the multiplier-one image.
The finite quotient $H_1=C/F$ is
\[
 \operatorname{Dic}_3=
 \langle r,s:r^6=1,\ s^2=r^3,\ srs^{-1}=r^{-1}\rangle,
\]
where $r,s$ are the images of $\rho,\sigma$.
If $u$ has multiplier one, it fixes $sr$ and sends
$r$ to a conjugate of $r$, and $s$ to a conjugate
of $s$. Hence
\[
         u(r)=r^\varepsilon,\qquad
         u(s)=r^{2k}s,\qquad
         2k-\varepsilon\equiv-1\pmod6.
\]
There are exactly two possibilities: the identity,
or $\operatorname{Int}(sr)$. Composing $u$ with the
inverse of $\operatorname{Int}(e_0^\varepsilon)$,
where now $\varepsilon\in\{0,1\}$ records this choice,
we may assume it is the identity on $H_1$.

On $M=F/F'=\widehat\mathbb Z x\oplus\widehat\mathbb Z y$,
conjugation by $\rho,\sigma$ has matrices
\[
 R=\begin{pmatrix}0&-1\\1&-1\end{pmatrix},\qquad
 S=\begin{pmatrix}0&1\\1&0\end{pmatrix}.
\]
The induced automorphism of $M$ commutes with both.
Solving these two matrix equations shows it is scalar;
because it fixes $x$, the scalar is one. The remaining
automorphism of $\Delta$ is thus a translation
one-cocycle $b_h\in M$ over the identity of $H_1$.
Write $A=b_\rho$, $B=b_\sigma$. The fixed loop gives
$B+SA=0$. On the other hand $\rho^3=\sigma^2$ and
$1+R+R^2=0$ give
\[
          0=(1+S)B=-(1+S)A.
\]
Therefore $A=B=t(1,-1)$ for a unique
$t\in\widehat\mathbb Z$. These are exactly
$(1-R)(tx)$ and $(1-S)(tx)$, so this automorphism
is $\operatorname{Int}(x^t)$. This proof uses the
integral equations directly; no division by two or
three is involved.

Since $x=\omega e_0^2$, the complete multiplier-one
image is $\operatorname{Int}(e_0^{\varepsilon+2t})$.
Every $r\in\widehat\mathbb Z$ has a unique such
decomposition after choosing its parity. The maps
$\operatorname{Int}(e_0^r)$ already come from $D$,
and (B.33.1) supplies every multiplier. This proves
surjectivity in (B.33.2). Faithfulness of $L_0$ gives
uniqueness of $\mu$ from the loop image, and
$L_0\cap Z(\Delta)=1$, proved in B.31, gives
uniqueness of $r$. Conjugating an inner loop power
by $s_\mu$ proves the stated group law. The two
generators of $\Delta$ give (B.33.3).

The executed certificate reduces this whole actual
image to $\Delta_n=C/(F'F^n)$ for $1\le n\le8$.
It verifies every element and every generator edge
for each displayed automorphism, and every pair in
the affine composition law. The image orders are
\[
\begin{array}{c|rrrrrrrr}
n&1&2&3&4&5&6&7&8\\\hline
|\Delta_n|&12&48&108&192&300&432&588&768\\
|\operatorname{im}D|&4&8&24&32&80&48&168&128.
\end{array}
\]
They equal $2n\varphi(2n)$ for even $n$, and
$2n\varphi(4n)$ for odd $n$. These are reductions
of the proven source image, not an identification
with every isolated signed solution on a finite group.
The input is the exact integral cocycle from B.24a;
its hash and all check counts are saved in the output.
\end{proof}

\begin{statement}{B.34}{Refuted}{The full profinite source does not act over the rational base}
There is no action of the whole actual source $D$ by
$\mathbb Q$-defined automorphisms of the minimal Lie
stage $\mathscr C_{0,\mathbb Q}$ which, after extension
to $\mathbb Q_5$, agrees with its actual continuous
action on $F^{\rm ab}\otimes_{\widehat\mathbb Z}\mathbb Q_5$.
Consequently no compatible action on an integral
$\mathbb Z$ model of this stage can have that property.
The continuous relative actions over $\mathbb Q_\ell$,
the abelian-stage actions over $\mathbb Z_\ell$, and
the profinite actions of B.31 are the valid replacement.
Actual discrete source automorphisms preserving the
integral lattices also act on the corresponding
discrete rational constructions.
\end{statement}
\begin{proof}
By Hensel's lemma there is a unique $a\in1+5\mathbb Z_5$
with $a^2=6$: start from the root one modulo five,
whose derivative is the unit two. This $a$ is not
rational, since a rational square equal to six would
have odd valuations at two and three. Set
\[
       \mu_5=a,\qquad\mu_p=1\quad(p\ne5).
\]
This is an element of $\widehat\mathbb Z^\times$.
The surjectivity used in B.33 supplies an actual
Galois, hence signed-source, element with multiplier
$\mu$. Its induced first-stage automorphism is
$s_\mu$. On the rank-two first Lie layer over
$\mathbb Q_5$ it is scalar multiplication by $a$.

A $\mathbb Q$-defined automorphism of the vector
group $\mathbb G_{a,\mathbb Q}^2$ is linear with a
rational matrix. Its trace remains rational after
extension to $\mathbb Q_5$, whereas the required
matrix has trace $2a\notin\mathbb Q$. This is the
contradiction, independent of any choice of rational
basis. An integral action would yield a rational
one on the generic fibre and is therefore impossible
as well. The program gives compatible Hensel roots
through $5^{12}$; the all-level existence and
irrationality are the preceding exact arguments.

The assertion concerns a Lie stage with its marked
continuous comparison. It does not deny the tautological
action on a constant profinite group scheme, and it
does not deny the rational first-stage existence
for rational fine data proved in B.24c. It specifies
why a profinite multiplier cannot automatically be
treated as a rational scalar. Formula (B.31.4) proves
the asserted local replacements. For a discrete
source automorphism the coset translations and
lattice maps are integral words; extending those
maps to the indicated rational constructions proves
the final assertion.
\end{proof}

\begin{statement}{B.35}{Proved}{Exact first-stage signed normalizers and the surviving midpoint character}
Use the actual sources of B.31, put $\Delta=C/F'$, and let
$A=\widehat{\mathbb Z}^{\times}$.  Write
\[
 \eta_8(\mu)=\frac{\mu_2^2-1}{8}\pmod2,
 \qquad
 \mathcal E_2=\{(t,\mu)\in\widehat{\mathbb Z}\rtimes A:
                         t\bmod2=\eta_8(\mu)\},
\]
where $(t,\mu)(v,\nu)=(t+\mu v,\mu\nu)$.
Then the images of the two actual signed normalizers in
$\operatorname{Aut}(\Delta)$ are
\[
\begin{aligned}
 \operatorname{im}N_\rho^*
   &=\{\operatorname{Int}(\rho^j)h_\mu:
                       j\in\mathbb Z/3,\ \mu\in A\}
     \simeq C_3\rtimes_{\epsilon_3}A,\\
 \operatorname{im}N_\sigma^*
   &=\{\operatorname{Int}((xy)^t\sigma^e)s_\mu:
              e\in\mathbb Z/2,\ (t,\mu)\in\mathcal E_2\}
     \simeq C_2\times\mathcal E_2.
                                                        \tag{B.35.1}
\end{aligned}
\]
Here the normalizer action is the specified strict action of B.31,
all displayed coordinates are unique, and
\[
 \begin{aligned}
 v_3(\mu)&=(1-\epsilon_3(\mu))/2,&
 a_\mu&=(\epsilon_3(\mu)\mu-1)/6,\\
 h_\mu&=\operatorname{Int}((xy^{-1})^{a_\mu}
                         \sigma^{v_3(\mu)})s_\mu.
 \end{aligned}
                                                        \tag{B.35.2}
\]
The quotient $\operatorname{im}N_\sigma^*\to A$ does not split.
In particular the retained $\sigma$ conjugator carries a genuine
additional coordinate: its multiplier-one translation image is
$2\widehat{\mathbb Z}(x+y)$, although the full quotient centralizer
contains $\widehat{\mathbb Z}(x+y)$.

The images of $G_D$, $\mathcal M_D$, and $S_0\mathcal M_D$,
using their projections in B.31, are all
\[
                  \operatorname{Inn}(\Delta)\rtimes s_A,
       \qquad Z(\Delta)=\langle\omega\rangle.             \tag{B.35.3}
\]
If $u|_\Delta=\operatorname{Int}(e_0^{r(u)})s_{\mu(u)}$,
the action kernel of $G_D$ is the concrete set
\[
 \{[c,u]:\mu(u)=1,\quad
                \overline c\,e_0^{r(u)}\in\langle\omega\rangle\}.
                                                        \tag{B.35.4}
\]
The two larger fibre-product kernels are its inverse images.
For the normalizers their kernels are the inverse images of
$(j,\mu)=(0,1)$ and $(e,t,\mu)=(0,0,1)$ in (B.35.1).
These conclusions concern the first modular stage $C/F'$;
they do not collapse the higher Fermat actions to this stage.
\end{statement}
\begin{proof}
We first identify the centralizers which are needed both for
normalization and for the geometric comparison.  Theorems 1 and 2
of \citet[printed p.~31]{B35HerfortRibes1984} state that torsion
elements of a free profinite product are conjugate into a factor,
and that the centralizer of a nontrivial factor element lies in
that factor.  The class of all finite groups satisfies their
closure hypotheses.  The opened primary text is their four-page
announcement with proof outline; we use its stated theorems,
not a discrete reduced-word argument applied to profinite words.
Applying the centralizer theorem to
$C/\langle\omega\rangle=C_3\widehat *C_2$ and lifting gives
\[
 C_C(\rho)=\langle\rho\rangle,
 \qquad C_C(\sigma)=\langle\sigma\rangle.                 \tag{B.35.5}
\]
The finite quotient $H_1=\operatorname{Dic}_3$ has corresponding
centralizers $\langle\overline\rho\rangle$ and
$\langle\overline\sigma\rangle$.  Strip these finite components
in $\Delta$ and use the matrices $R,S$ of B.33.  Since
$M^R=0$ and $M^S=\widehat{\mathbb Z}(x+y)$, this proves
\[
 C_\Delta(\rho)=\langle\rho\rangle,
 \qquad
 C_\Delta(\sigma)=\langle\sigma\rangle
                         \times\widehat{\mathbb Z}(x+y).
                                                        \tag{B.35.6}
\]
The kernel of $H_1\to\operatorname{Aut}(M)$ is its central
$\langle\omega\rangle$, and $M^{H_1}=0$.  Thus a central
element of $\Delta$ is exactly $1$ or $\omega$.

By B.33 the image of $G_D$ is generated by
$\operatorname{Inn}(\Delta)$ and $s_A$.  These intersect trivially:
on $M$, an inner automorphism acts by an element of the displayed
$S_3$ representation, while $s_\mu$ acts by the scalar $\mu$;
the only scalar in this integral $S_3$ representation is the identity.
Thus $\mu=1$ and $s_\mu=1$.  This proves the semidirect product
in (B.35.3), and directly gives (B.35.4).  Both normalizers map
onto $\gamma$: for a representative $u\in D$ of a class in
$\gamma$, a conjugator for $u(\rho)$, respectively $u(\sigma)$,
corrects $u$ to an element of the relevant normalizer.
Consequently the projections $\mathcal M_D\to G_D$ and
$S_0\mathcal M_D\to G_D$ are onto.  This proves their image
and kernel assertions without replacing the conjugator source
by a group of unmarked automorphisms.

\emph{The $\rho$ image.}
Put $n=(\mu-1)/2$.  Direct calculation in $\Delta$ gives
\[
 s_\mu(\rho)=y^n\rho,\qquad
 s_\mu(\sigma)=\sigma^{\epsilon_2(\mu)},\qquad
 \sigma\rho\sigma^{-1}=x\rho^{-1}.                       \tag{B.35.7}
\]
For example use $\sigma=e_0e_1e_0$ and
$\rho=e_0^{-1}e_1^{-1}$, together with
$e_0^2=\omega x$, $e_1^2=\omega y$; the central sign in
$s_\mu(\sigma)$ is fixed by its abelianization modulo four.
If $\epsilon_3(\mu)=1$, conjugation by
$(xy^{-1})^{n/3}$ removes the translation in $s_\mu(\rho)$.
If $\epsilon_3(\mu)=-1$, first conjugate by $\sigma$ and then
by $(xy^{-1})^{-(n+1)/3}$.  The divisions are integral because
$\mu$ is odd and has its specified residue modulo three.
This proves $h_\mu(\rho)=\rho^{\epsilon_3(\mu)}$.
Moreover
\[
 a_\mu+\epsilon_3(\mu)\mu a_\nu=a_{\mu\nu}.
\]
The action of $\sigma$ on $x-y$ is $-1$, and its square is
central.  This identity therefore proves
$h_\mu h_\nu=h_{\mu\nu}$ as automorphisms of $\Delta$.
Conjugation by $h_\mu$ acts on $\operatorname{Int}(\rho)$
through $\epsilon_3(\mu)$.

Every multiplier occurs in $N_\rho^*$, by its surjection to
$\gamma$ and B.33.  Any two target normalizers with the same
multiplier differ by $\operatorname{Int}(C_\Delta(\rho))$,
which is exactly $C_3$ by (B.35.6).  This entire finite kernel
already comes from $\langle\rho\rangle\subset C$.
Thus every $h_\mu$ and exactly the stated three elements above
it occur.  This proves the first line of (B.35.1).

\emph{An integral restriction on the $\sigma$ coordinate.}
Equation (B.35.7) and (B.35.6) show that every ambient target
normalizer has a unique expression
$\operatorname{Int}((xy)^t\sigma^e)s_\mu$, with
$t\in\widehat{\mathbb Z}$ and $e\in\mathbb Z/2$.
We determine which are supplied by the actual source.
First suppose its multiplier is one.  Write the source element
as $[c,u]$ and absorb the inner loop factor of B.33 into $c$.
We may therefore assume $u\in D_0$, so that $u$ is the identity
on $\Delta$.  There are unique normalized conjugators
$\alpha,\beta\in F$ satisfying
\[
          u(\rho)=\alpha\rho\alpha^{-1},\qquad
          u(\sigma)=\beta\sigma\beta^{-1}.               \tag{B.35.8}
\]
Existence follows by reducing a conjugator in $H_1$ and
multiplying it by a power of the corresponding finite-order
generator; uniqueness follows from (B.35.5) and
$F\cap\langle\rho\rangle=F\cap\langle\sigma\rangle=1$.
The equations modulo $F'$ give
\[
             \alpha^{\rm ab}=0,\qquad
             \beta^{\rm ab}=k(x+y).
\]
Here $1-R$ is injective over $\widehat{\mathbb Z}$, even
though its determinant three is not a unit at three.

The action of $R$ on $\gamma_2F/\gamma_3F$ is its determinant,
namely one.  Since $\alpha\in F'$, (B.35.8) implies
$u(\rho)\rho^{-1}\in\gamma_3F$.  The fixed loop $\sigma\rho$
then gives $u(\sigma)\sigma^{-1}\in\gamma_3F$ as well.
In the integral class-two quotient write uniquely
\[
             \beta=x^k y^k c^q,\qquad c=[x,y].
\]
The interchange sends this element to
$x^k y^k c^{-k^2-q}$.  Thus
\[
                        2q+k^2=0.
                                                        \tag{B.35.9}
\]
Reduction modulo two forces $k\in2\widehat{\mathbb Z}$.
No division by two in a Malcev algebra has been used.
Finally the normalizer equation says $c\beta\in C_C(\sigma)$,
so its target translation is $-k$.  The actual multiplier-one
image is therefore contained in
$C_2\times2\widehat{\mathbb Z}(x+y)$.

\emph{Midpoint realization in the actual arithmetic subgroup.}
Here we justify the geometric comparison and its basepoint
ambiguity explicitly.  Let $U=\mathbb P^1-\{0,1,\infty\}$,
with tangent $0+1$.  The quotient stack $[U/S_3]$ has orbifold
fundamental group $C_3\widehat *C_2=C/\langle\omega\rangle$:
the invariant map
\[
               j(z)=\frac{256(1-z+z^2)^3}{z^2(1-z)^2}
\]
has stabilizers of orders three and two over $0$ and $1728$,
and its cusp is removed.  Its index-six kernel is $\pi_1(U)$.
In the braid presentation, the two standard meridians are
$a^2,b^2$, whose images are $x^{-1},y^{-1}$.
The standard arithmetic GT action of E.6 restricts to the usual
based Galois action on these meridians.  It is consequently the
arithmetic action on this quotient stack.  Indeed an automorphism
of $C/\langle\omega\rangle$ which fixes $F$ pointwise is the
identity: its discrepancy on any element centralizes $F$.
That centralizer is trivial.  To see this, a centralizing element
has a power in $F$, hence a power in $Z(F)=1$; the torsion and
centralizer theorems quoted above then exclude a nontrivial
finite-order element centralizing the nonabelian open subgroup.
Thus the extension of the specified $F$ action is unique.

The relevant involution of $U$ is $h(z)=1-z$, with rational
fixed point $P=1/2$.  The subgroup above $\langle h\rangle$ is
\[
 F\rtimes\langle\overline\sigma\rangle
     =\langle x,\overline\sigma\rangle
     \simeq\widehat{\mathbb Z}\widehat *C_2:
\]
eliminate $y=\overline\sigma x\overline\sigma^{-1}$ in its
profinite presentation.  The stabilizer at $P$ is an order-two
subgroup of this group.  Herfort--Ribes Theorem 1 conjugates
it to the $C_2$ factor, since $\widehat{\mathbb Z}$ is
torsion-free.  The conjugator can be taken in $F$: multiplying
by its $C_2$ component does not change the conjugated factor.
Choose an etale path from the tangent to $P$ realizing that
conjugacy.  If $g\in G_\mathbb Q$, its path cocycle is a
conjugator $\beta_g$ for the involution, since the stabilizer
at the rational fixed point is constant downstairs.
Upstairs in $C$ the possible extra $\omega$ factor is excluded
by the forced image of $\sigma$ in $C^{\rm ab}$; the exponent
is precisely $\epsilon_2(\chi(g))$.  Hence
$[\beta_g^{-1},g]\in N_\sigma^*$ for the actual arithmetic
source, with the signs of B.31 retained.

The covers obtained by adjoining roots of $z$ and $1-z$
compute the two abelian coordinates of this path cocycle.
Their values at $1/2$ are both the Kummer class of $1/2$;
this is also \citet[Lemma 3.16]{FA_Wickelgren2012}.
The same direct cover calculation is valid at every prime,
including two.  Reversing the path or the two meridian
orientations changes the common sign.  Changing the comparison
path adds $(\chi(g)-1)v$ to its abelian cocycle for some
$v\in\widehat{\mathbb Z}^2$.  Consequently, on
$G_{\mathbb Q(\mu_\infty)}$ its diagonal coordinate is
exactly $\pm\kappa_2(g)$, and for arbitrary $g$ its parity
is that of $\kappa_2(g)$, independently of either convention.
These are the only two assertions needed for the exact image.
The diagonal equality itself also follows from
$g|_\Delta=s_{\chi(g)}$ and (B.35.7).

For completeness the image of this Kummer character on the
cyclotomic kernel is exactly
\[
          \kappa_2(G_{\mathbb Q(\mu_\infty)})
                              =2\widehat{\mathbb Z}.      \tag{B.35.10}
\]
The image is a closed subgroup of $\widehat{\mathbb Z}$.
Its divisibility by $p^r$ is equivalent to
$2^{1/p^r}\in\mathbb Q(\mu_\infty)$.
The square root belongs to $\mathbb Q(\zeta_8)$.
For every integer $m>2$, $X^m-2$ is Eisenstein, and
$\mathbb Q(2^{1/m})$ is real but has nonreal conjugates;
it is therefore not normal over $\mathbb Q$ and cannot
be a finite subextension of the abelian extension
$\mathbb Q(\mu_\infty)/\mathbb Q$.  Thus the two-primary
image has valuation exactly one and every other primary
image has valuation zero, proving (B.35.10).
Finally
$\sqrt2=\zeta_8+\zeta_8^{-1}$ gives
\[
        \kappa_2(g)\pmod2=\eta_8(\chi(g)).                \tag{B.35.11}
\]
Every multiplier occurs arithmetically.  Equations
(B.35.10)--(B.35.11) supply every translation in the claimed
parity coset; (B.35.9), applied to the quotient of any two
actual normalizers having the same multiplier, excludes all
others.  The finite $C_2$ factor already comes from $\sigma$.
This proves the second line of (B.35.1).  Its law is the
stated affine law, because $s_\mu$ scales $x+y$ by $\mu$;
$\operatorname{Int}(\sigma)$ is central in this normalizer
image, since inversion has no effect on its order-two class.

\emph{Nonsplitting and finite certificate.}
A section over $A$ would give a one-cocycle $t:A\to\widehat{\mathbb Z}$
with $t(\mu)\bmod2=\eta_8(\mu)$.  Since $A$ is abelian,
commuting $\mu$ with $-1$ gives
\[
                  2t(\mu)=(1-\mu)t(-1).
\]
But $\eta_8(-1)=0$, so $t(-1)=2b$ and
$t(\mu)=(1-\mu)b$ is always even.  This contradicts any
$\mu$ with $\mu_2\equiv3\pmod8$.  Projection to
$\mathcal E_2$ shows that the extra central $C_2$ cannot
remove the obstruction.

The program \texttt{g\_signed\_normalizers.py} reconstructs
$\Delta_n=C/(F'F^n)$ from the exact integral cocycle of B.24a,
for $1\le n\le8$. It checks (B.35.7), the section law in
(B.35.2), and the reductions of both normalizer images on
all group elements, with all generator-edge homomorphism
checks. It also checks (B.35.9) modulo $2,4,\ldots,64$.
Nonsplitting is already visible with translation modulo four
and multiplier modulo eight: none of the four pairs of lifts
of $3,-1\in(\mathbb Z/8)^\times$ commute.
These are finite checks of the proved source images;
arithmetic realization is supplied by the preceding path
and Kummer argument, not inferred from finite shadows.
\end{proof}

